% ==========================================================================
% GENERATED by scripts/split_paper.py. Do not edit this file.
%
% Source: imported_papers/sds_hounsfield_features/main.tex
% Part 3: Hounsfield Unit Ranges as Inductive Bias
%
% Edits belong upstream, in the published paper. To change what the
% chapter includes, change imports/hounsfield.toml and re-run
% the splitter.
% ==========================================================================
\section{Hounsfield Unit Ranges as Inductive Bias}
The goal of tissue segmentation for BCA is to classify each pixel of a 2D CT slice into a tissue class or the background class.
We denote the CT slice, a 2D image of HU values, as matrix $\bm{X}_0 \in \mathbb{R}^{h, w}$, where $h$ and $w$ denote the height and width of the slice.
The HU pixel values in $\bm{X}_0$ typically range from $-\num{1000}$ (air) to approximately $+\num{3000}$ (dense bone) \cite{bibb2014medical}, while foreign objects such as metal implants can exhibit HU values up to $+\num{30000}$.
The pixel-wise classification outputs a segmentation mask $\hat{\bm{Y}} \in (\mathcal{C} \cup \{ BG \})^{h , w}$ of the same shape as the input image, where $\mathcal{C}$ represents the set of tissue classes and $BG$ the background class.

\Cref{fig:overview} visualizes how we integrate the prior knowledge about tissue-specific HU ranges into the model.
On top of the input $\bm{X}_0$ we stack $|\mathcal{C}|$ binary masks $\{ \bm{X}_c\ |\ c \in \mathcal{C}\}$, yielding $\bm{X} \in \mathbb{R}^{(1 + |\mathcal{C}|) , h , w}$. 
Given $\bm{X}_0(i,j)$ be the pixel value at row $i$ and column $j$ in matrix $\bm{X}_0$, $\bm{X}_c$ for class $c$ is defined as:
% 
\begin{equation*}
\bm{X}_{c}(i,j)=\begin{cases}
    1, & \text{if } \bm{X}_{0}(i,j) \text{ within the HU range of class } c,\\
    0, & \text{otherwise}.
\end{cases}
\label{eq:pixel_indicator}
\end{equation*}

Thus, the first channel of the input $\bm{X}$ contains the ``original'' CT image $\bm{X}_0$, while the additional channels contain the binary masks $\bm{X}_c$, each indicating if a class $c$ is plausible at specific pixel locations according to medical textbook knowledge on HU ranges.
%
\begin{table}[t]
    \caption{HU ranges for each tissue type.}
    \begin{center}
        \begin{tabular}{lcc}
        \toprule
        \textbf{Tissue type} & \multicolumn{2}{c}{\textbf{HU range}} \\
        \cmidrule(lr){2-3}
         & \textbf{From} & \textbf{To} \\
        \midrule
        Skeletal muscle (SM) & -29 & 150 \\
        Intramuscular adipose tissue (IMAT) & -190 & -30 \\
        Visceral adipose tissue (VAT) & -150 & -50 \\
        Subcutaneous adipose tissue (SAT) & -190 & -30 \\
        \bottomrule
        \end{tabular}
    \label{tab:hu_values}
    \end{center}
\end{table}
%
\Cref{tab:hu_values} provides an overview of the tissue types and corresponding HU ranges used in this study. 
Equivalent HU ranges (IMAT, SAT) lead to identical binary masks and could also be represented as a single mask.
Overlapping HU ranges (e.g., VAT, SAT) lead to overlapping masks, indicating multiple possible tissue types at some locations.
Adding the respective information as additional input channels effectively constrains the model’s decision space, reducing ambiguity and improving segmentation accuracy by leveraging well-established radiodensity properties.
